\cotefile{157-5}
\scanpage{101}
\mspage{\unsure{88|98}}
% archive sheet number: 100

\uline{Dém} La nécessité de l'axiome de factorisation
pour $(TC,F)$, pour avoir un \unsure{triple} de Quillen,
est tautologique. Pour la suffisance, ayant
déjà les deux \unsure{principes|couples|quintuples} de Quillen $(C,TF),(TC,F)$
avec
\diagram{\begin{tikzcd}[ampersand replacement=\&,row sep=small,column sep=large]
TC \arrow[d,leftrightarrow] \&[-1em] \subset \&[-1em] C \arrow[d,leftrightarrow]\\
F \& \unsure{\supset|\longrightarrow} \& TF
\end{tikzcd}}
, et $TC=C\cap W$, le fait que \deleted{\unsure{c'est une condition}} \inserted{$(W,C,F)$ est}
un triple de Quillen est le cor.~4 p.~35
de $\overline{\underline{\mathrm{XV}}}$) : les conditions $TC=C\cap W$, $TF=F\cap W$,
$W=TF\circ TC$ sont équivalentes, et équivalent au
fait que $(W,C,F)$ soit un triple de Quillen.

Par la définition de Quillen des triples \uline{clos},
\unsure{par|vu} les propriétés
\[
\underbrace{W=TF\circ TC}_{\substack{\text{vraie pour tous}\\\text{les triples}\\\text{de Quillen}}},\qquad
\deleted{\unsure{E|F}}C=(TF)^{*},\qquad
\underbrace{F=(TC)_{\unsure{*|x}}}_{\substack{\text{vraie dans}\\\text{le cas actuel}}},
\]
il reste à exprimer que $C=(TF)^{*}$, ce qui
équivaut à : $C$ stable par facteurs directs\inserted{(cf loc. cit. cor.~1 p.~31)},
qed.

\uline{Remarques} 5.4. Les \deleted{\unsure{prop.}} énoncés 5.2. et 5.3,
et la relation $T\widetilde C\subset (TC)^{\unsure{4|\natural}}$ \unsure{aux|au} prop.~5.1,
ne font pas appel à la totalité de l'hypothèse que $(M,W,C)$ soit une catégorie à cofibrations spéciale, mais seulement que l'on ait
une cat. à cofibrations avec
\diagram{\begin{tikzcd}[ampersand replacement=\&,column sep=small,row sep=tiny]
C \arrow[r,leftrightarrow,"\shortmid" description] \& TF \arrow[d,equals,"\mathrm{déf}"] \&[-1em] \subset W\\
\& C_{*} \&
\end{tikzcd}}
\unsure{;|,}

\scanpage{102}
\mspage{\unsure{89|109}}
% archive sheet number: 101

i.e. une catégorie avec $W,C$ \deleted{\unsure{fibr. stable à}, $C$ stable par} \inserted{satisfaisant}
$\mathrm{Cof}_{0},\mathrm{Cof}_{1},\mathrm{Cof}_{2}$ de la page 39, \unsure{plus} $TF\overset{\mathrm{déf}}=C_{*}\subset W$,
plus l'axiome de factorisation pour $(C,TF)$
(lequel implique $\mathrm{Cof}_{3}$).

\uline{Corollaire} 5.5. Soit \deleted{\unsure{$C$,}} \inserted{$(M,W,C)$} une catégorie à cofibrations
spéciale, telle qu'il existe une \uline{petite}
partie $TC_{0}\subset TC=C\cap W$, avec
\[
TC\subset {TC_{0}}^{\sim}, \tag*{(5.5.1)}
\]
ou ce qui revient au $=$, $(TC)^{\sim}=(TC_{0})^{\sim}$ i.e.
\[
\underset{\substack{\Vert\,\mathrm{déf}\\(TC)_{*}}}{F}=(TC_{0})_{*}.\tag*{(5.5.2)}
\]
Alors \deleted{\unsure{$(W,C$|$(M,W,C$}} $(W,C,F)$ est un triple de
Quillen. \deleted{clos ssi} Ce triple est clos ssi $C$ est
stable par facteurs directs.

En effet, le th. de factorisation de
\uline{XV} \S~1 \unsure{implique|infère} que moyennant l'hypo \unsure{faite|faible}
\deleted{\unsure{sur}} $TC_{0}\subset TC$, le couple $(TC,F)$ \deleted{\unsure{satisfait}}
satisfait à la condition de factorisation.
En effet, par les conditions $b,b'$ de la
définition 4.4., $TC$ a les propriétés de
stabilité voulues (\unsure{contient} $\unsure{\mathrm{iz}|\mathrm{is}}$, stable par
\unsure{cochgt|cochgt.} de base et par composition transfinie).

\scanpage{103}
\mspage{\unsure{90}}
% archive sheet number: 102

\deleted{\uline{\unsure{Définition}}} \inserted{\uline{Définition} 5.6.} Les triples de Quillen \unsure{déduits} \inserted{comme dans 5.5} d'une
catégorie à cofibrations spéciale sont dits
\uline{spéciaux}. Ainsi un \unsure{triples|triple} de Quillen
$(W,C,F)$ de $M$ est spécial s'il
satisfait les conditions a) à d) de la
définition (4.4.), concernant $(W,C)$ exclusivement, plus les deux conditions suivantes

\textcircled{\scriptsize $c'$} $\exists$ petite partie $TC_{0}$ de $TC$, telle
que $TC\subset (TC_{0})^{\sim}$.

\[
F\overset{\deleted{\unsure{\mathrm{déf}}}}{=}(TC)_{*}. \tag*{\textcircled{\scriptsize $\unsure{e|e'}$}}
\]

\uline{Corollaire} 5.7. Soit \deleted{\unsure{$M$,}} $(W,C,F)$ un triple

\noindent\vrule\hspace{.6em}\begin{minipage}{.93\linewidth}
de Quillen \uline{clos} dans $M$. Pour que
ce triple soit spécial, il f et s.
qu'il satisfasse les conditions suivantes

\textcircled{\scriptsize $a$} $M$ stable par \inserted{petites} $\unsure{\varprojlim|\varinjlim}$,
et ses éléments sont accessibles

\textcircled{\scriptsize $\unsure{c|b},c'$} Il existe une petite partie $C_{0}$ de $C$,
$TC_{0}$ de $TC$, telles que l'on ait
\[
 C\deleted{\unsure{\subset}}={C_{0}}^{\sim},\qquad TC=(TC_{0})^{\sim}.\tag*{$(5.7.1)_{a}$}
\]
\end{minipage}

En effet, toutes les \inserted{autres} conditions \deleted{\unsure{sont}}
a) à e), y inclus c') de 5.6., sont
automatiquement satisfaites.

\scanpage{104}
\mspage{91}
% archive sheet number: 103

\uline{Proposition} 5.8. Considérons un couple de
foncteurs adjoints
\diagram{\begin{tikzcd}[ampersand replacement=\&,column sep=large]
{(5.8.1)} \& M \arrow[r,dash,shift left=1.1ex,"{u_{!}}"] \& M' \arrow[l,shift left=1.1ex,"{u^{*}}"]
\end{tikzcd}}
, où $M$ est muni d'un triple de Quillen
spécial \inserted{clos} $(W,C,F=(TC)_{*})$, avec $TC=C\cap W$
(cf définition 5.6.). \deleted{Définissons \unsure{$W',C',F'$ de $M'$|$W',C',F'$}} \inserted{On suppose} \deleted{\unsure{que}} \deleted{comme ds prop.~4.6. (cf (4.6.1))} \deleted{$M'$ est stable par petites $\varinjlim$, que} \unsure{sont}
satisfaites les conditions suivantes (comparer
prop.~4.6., relations (4.\unsure{8|5|6}.1), et (4.7.4)) :
\[
\left\{\begin{aligned}
 &a)\ M'\ \text{stable par petites }\varinjlim,\ \text{ses éléments}\\
 &\hspace{2.2em}\text{sont accessibles}\\
 &b)\ u^{*}\ \text{commute aux petites }\varinjlim\\
 &c)\ u^{*}(u_{!}(C))\subset C,\qquad u^{*}(u_{!}(TC))\subset TC.
\end{aligned}\right. \tag*{(5.8.2)}
\]
Définissons $(W',C',F')$ dans $M'$ (comparer
prop.~4.6, \unsure{notamment|notations|relations} 4.6.2) par
\[
\left\{\begin{aligned}
 W'&=(u^{*})^{-1}(W)\\
 C'&=\deleted{\text{clôture de }}u_{!}(C)^{\sim}\deleted{\text{ par cochgt. de}}\\
 &\quad\deleted{\text{base et par composition transfinie}}\\
 F'&=(TC')_{*},\qquad\text{où }\ TC'\overset{\mathrm{déf}}=C'\cap W'.
\end{aligned}\right. \tag*{(5.8.3)}
\]
Alors $(W',C',F')$ est un triple de Quillen
clos spécial dans $M'$, et on a de plus
\[
 F'=(u^{*})^{-1}(F),\qquad TF'=(u^{*})^{-1}(TF)\tag*{(5.8.4.)}
\]
\[
 TC'=(u_{!}(TC))^{\sim}.\tag*{(5.8.5)}
\]

\scanpage{105}
\mspage{92}
% archive sheet number: 104

\uline{Démonstration} Par 4.6, on sait déjà que
$(M',W',C')$ est une cat. à cofibrations spéciale.
[\uline{NB} Ce qui est noté $C'$ ici, était noté $\widetilde{C'}$ dans
la prop.~4.5. On doit utiliser la dernière
assertion de 4.6, où on suppose $C$ stable par
facteurs directs, OK car $(W,C,F)$ supposé
clos] Pour prouver que $(W',C',F'\overset{\mathrm{déf}}=(TC')_{*})$
est un triple de Quillen spécial, il
reste, par 5.3. et définition 5.6, à prouver
qu'il existe une petite partie $TC'_{0}$ de
$TC'_{\unsure{9|0}}$, telle que
\[
 (TC'_{0}\unsure{\subset|\subseteq})\quad TC'\subset (TC'_{0})^{\sim}. \tag*{(5.8.6)}
\]
(Le triple de Quillen sera de plus clos, puisque
$\deleted{\unsure{T}}C'$ est \inserted{\unsure{Q-}}clos \inserted{\unsure{à g.|à g}} par définition.) Soit donc
$TC_{0}$ une petite partie de $TC$ telle que
\[
 TC\subset (TC_{0})^{\sim},\tag*{(5.8.7)}
\]
qui existe par hypothèse sur $(W,C,F)$. On
prendra
\[
 TC'_{0}\overset{\mathrm{déf}}=u_{!}(TC_{0}).\tag*{(5.8.8)}
\]
On a vu dans 4.7. qu'on a
\[
 u_{!}(TC)\subset TC'\tag*{(5.8.9)}
\]
(\unsure{résulte} de l'hyp. $u^{*}(u_{!}(TC))\subset TC$
de (5.8.2 c)), \unsure{a|à} fortiori
\[
 TC'_{0}=u_{!}(TC_{0})\subset TC',
\]

\scanpage{106}
\mspage{93}
% archive sheet number: 105

il reste à prouver
\[
 TC'\overset{?}{\subset}(TC'_{0}\unsure{{}'})^{\sim}\tag*{(5.8.10)}
\]
ce qui équivaut à $(TC')^{\sim}\subset (TC'_{0})^{\sim}$ ou
encore à
\[
 (TC'_{0})_{*}\overset{?}{\subset}(TC')_{*}\overset{\mathrm{déf}}=F',\tag*{(5.8.\unsure{10 bis})}
\]
alors qu'on a a priori l'inclusion
inverse (\unsure{qui suit|qui vient|qui résulte} de $TC'_{0}\subset TC'$)
\[
 F'\subset (TC'_{0})_{*}.\tag*{(5.8.11)}
\]
Or par adjonction on a
\[
\begin{aligned}
 (\underset{u_{!}(TC_{0})}{\underline{TC'_{0}}})_{*}
 &=\deleted{\unsure{(u_{!}}}(u_{!}(TC_{0}))_{*}=(u^{*})^{-1}\bigl(\underset{\substack{\unsure{\Vert|\mid F}\\\unsure{*}(TC)_{*}\overset{\mathrm{déf}}=F}}{\underline{(TC_{0})_{*}}}\bigr)\\
 &=\unsure{V|u}^{*\, -1}(F),
\end{aligned}
\]
donc (5.8.11.) s'écrit
\[
 F'\subset u^{*\, -1}(F)\qquad\text{i.e.}\qquad u^{\unsure{4|*}}(F')\subset F,\tag*{(5.8.12)}
\]
Comparer (4.7.6)), et il faut montrer
que cette inclusion est une égalité
(sur quoi j'avais buté p.~100, 101\unsure{\ldots|--}),
ce qui n'est autre que la première des
deux relations (5.8.4.). \deleted{\unsure{Celle-ci}} \deleted{Mais je vais}
\deleted{\unsure{indiquer dé}} Mais c'est sous la forme \deleted{\unsure{plus spéciale}} \inserted{\unsure{initiale|inutile}}
(5.8.10) \inserted{(\unsure{bien près de|donc aussi de|bien pire de} (5.8.4))} que nous allons prouver la chose,
finalement\unsure{.|,} \deleted{\illegible{1}} Notons que l'on a

\scanpage{107}
\mspage{94}
% archive sheet number: 106

\[
 \bigl(\underset{TC'_{0}}{\underline{u_{!}(TC_{0})}}\bigr)^{\sim}=(u_{!}(TC))^{\sim},
\]
puisque $TC$ est la clôture de $TC_{0}$ pour
les opérations de cochgt. de base, \deleted{\unsure{et}} de
composition transfinie et de facteurs directs,
et que $u_{!}$ commute à ces opérations. Donc
on doit montrer finalement
\[
 TC'\subset (u_{!}(TC))^{\sim}\tag*{(5.8.12 bis)}
\]
(où il n'est plus question de $TC_{0}$).
Mais le th. de factorisation s'applique
au couple
\[
 (u_{!}(TC))^{\sim},\ (u_{!}(TC))_{*}=(u^{*})^{-1}(F)\big), 
\]
puisque $(u_{!}(TC))^{\sim}=(TC'_{0})^{\sim}$, $TC'_{0}$ petit.
Si donc on a $f'\in TC'=C'\cap W'$,
\unsure{factorisant|factorisons} $f'$ en $p'i'$,
\diagram{\begin{tikzcd}[ampersand replacement=\&,column sep=small,row sep=large,execute at end picture={\draw (current bounding box.north west) -- (current bounding box.north east) -- (current bounding box.south east) -- (current bounding box.south west); \node[anchor=east] at ([xshift=-.8em]current bounding box.west) {$(\alpha)$};}]
X' \arrow[rr,dash,"{f'\in C'\cap W'}"] \arrow[dr,"{i'\,\unsure{\in|\notin}}"'] \& \& Y' \arrow[dl,dashed,shift right=1.2ex,"{q'}"']\\
\& X'_{1} \arrow[ur,"{p'}"'] \&
\end{tikzcd}}
avec \deleted{$i'\in (u_{!}(TC))^{\sim}$, \unsure{$p'\in$} $p'\in(u^{*})^{-1}(F)$, \unsure{Noto}}
\[
 i'\in (u_{!}(TC))^{\sim},\qquad p'\in (u^{*})^{-1}(F).\tag*{(5.8.13.)}
\]
Notons qu'on a \deleted{(\unsure{puisque}) à cause de (5.8.9)}
\deleted{\[
 (u_{!}(TC))^{\sim}\subset \deleted{\unsure{C'}}(TC')^{\sim}
\]}
\deleted{et d'autre part je dis qu'on a}
\deleted{\[
 (TC')^{\sim}\subset TC'\deleted{\unsure{\cap W'}}\ (=\unsure{C|C'}\cap W')\tag*{(5.8.14)}
\]}
\deleted{i.e. $TC'$ est Q-clos à \unsure{g.|G.|dr.} En effet, comme \unsure{on a} d'une part $(\underset{C'\cap W'}{TC'})^{\sim}\subset {C'}^{\sim}=C'$, de sorte qu'il suffit de voir que}

\scanpage{108}
\mspage{95}
% archive sheet number: 107

\deleted{\[
 (TC')^{\sim}\subset W'\qquad\text{i.e.}\qquad u^{*}((TC')^{\sim})\subset W'
\]}
\deleted{$C'$ est Q-clos, pour prouver $(TC')^{\sim}\subset TC'\unsure{\mathrel{\not=}|(}\overset{\mathrm{déf}}=C'\cap W'$, il suffit de prouver $(TC')^{\sim}\subset W'$ i.e.}
\deleted{\[
 u^{*}(TC')\subset W
\]}
\[
 (u_{!}(TC))^{\sim}\subset TC'\,(\overset{\mathrm{déf}}=C'\cap W')\tag*{(5.8.14)}
\]
car $u_{!}(TC)\subset TC'\subset C'$ donc $(u_{!}(TC))^{\sim}\subset {C'}^{\sim}=C'$,
et il reste à prouver que $u_{!}(TC)^{\sim}\subset W'$,
i.e. $u^{*}((u_{!}(TC))^{\sim})\subset W$. Or je dis qu'on
a même
\[
 u^{*}((u_{!}(TC))^{\sim})\subset TC \tag*{\deleted{\unsure{(5.8.15)}}}
\]
i.e.
\[
 u_{!}(TC)^{\sim}\subset (u^{*})^{-1}(TC).\tag*{(5.8.15)}
\]
Cela résulte du fait que $TC$ est Q-clos
donc Q-saturé à gauche, ce qui implique,
comme $u^{*}$ commute aux \inserted{petites} $\unsure{\varinjlim|\underline{\lim}}$, que
$u^{*\, -1}(TC)$ l'est également. D'autre part,
il contient $u_{!}(TC)$, donc aussi le
Q-saturé à gauche de $u_{!}(TC)$. Mais
ce dernier est \uline{égal} à la Q-clôture
$(u_{!}(TC))^{\sim}$, par le th. de factorisation,
vu que $\deleted{\unsure{u_{!}}}(TC'_{0})\subset u_{!}(TC\deleted{\unsure{{}_{0}}})\subset\deleted{\unsure{(u_{!}}}(TC'_{0})^{\sim}\deleted{\unsure{)}}$,
avec \deleted{\unsure{$u_{!}(TC_{0})$}} $TC'_{0}$ petit.

Ainsi, (5.8.13) et (5.\deleted{\unsure{4}}8.14) impliquent
$i'\in TC'=C'\cap W'$, donc $i'\in W'$, d'où $p'\in W'$,
i.e. \deleted{$p'\in$} $u^{*}(p')\in W$, et comme par hypo

\scanpage{109}
\mspage{96}
% archive sheet number: 108

(5.8.13) on a $u^{*}(p')\in F$, on a $u^{*}(p')\in TF$,
i.e.
\[
 p'\in(u^{*})^{-1}(TF)=TF'\unsure{6|,}(\overset{\mathrm{déf}}=C'_{*}).
\]
la \deleted{dernière} égalité \deleted{\unsure{étant}} $(u^{*})^{-1}(\underset{\deleted{C_{*}}}{\underline{TF}})=TF'$,
\deleted{résultant de l'adjonction}
i.e.
\[
 (u^{*})^{-1}(C_{*})=(u_{!}(C))_{*}\ \unsure{\times|,}
\]
résultant de l'adjonction de $u_{!},u^{*}$. Au
total, on a donc
\[
 p'\in TF'\unsure{\times|,} (=C'_{*})\tag*{(5.8.16)}
\]
et comme $f'\in C'$, le diagramme $(\alpha)$
p.~114 nous montre qu'il existe
\[
 q':Y'\mathrel{\unsure{-|\longrightarrow}}X'_{1}\unsure{.,|,}
\]
avec $p'q'=\mathrm{id}$, $q'f'=i'$. Donc $f'$ est
facteur direct de $i'\in u_{!}(TC)^{\sim}$, donc
on a $f'\in u_{!}(TC)^{\sim}$, d'où la relation
\[
 TC'\deleted{\unsure{\cap W'}}\subset u_{!}(TC)^{\sim},
\]
ce qui achève de prouver que $(W',C',F')$
est un triple de Quillen spécial et clos.

Les deux relations (5.8.4) ont été prouvées
chemin faisant (la deuxième étant essentiellement triviale, et la première représentant
la difficulté principale \unsure{dans|de} la démonstration\unsure{.|)}
Reste \unsure{à établir|'établir} (5.8.5.). Mais on sait
déjà \deleted{\unsure{que|par}}(5.8.9) que $u_{!}(TC)\subset TC'$,

\scanpage{110}
\mspage{97}
% archive sheet number: 109

d'où
\[
 u_{!}(TC)^{\sim}\subset TC'
\]
puisque $TC'$ est Q-clos à gauche, \deleted{\unsure{(ceci}} \inserted{(ou parce
que ça a été prouvé dans (5.8.14))}, \deleted{résulte d'autre part} on a \unsure{déjà} prouvé l'inclusion en sens inverse (5.8.12 bis),
\unsure{ess.} équivalente à 5.8.4 (et qui représentait
la principale difficulté).

\uline{Remarque} 5.9. Si on ne suppose pas que la
\deleted{soit} structure sur $M$ soit close, alors
la démonstration donnée de 5.8. montre
qu'on trouve \deleted{\unsure{quand}} sur $M'$ une
structure de Quillen spéciale, en définissant
cette fois $C'$ comme dans 4.6. Mais je
doute qu'on aura \inserted{jamais} à travailler avec des
triples \unsure{de} Quillen spéciaux qui ne soient
clos. De toute façon, si $W$ est
stable par petites $\varinjlim$ filtrantes\textsuperscript{$*$},
% slanted
\margin{ce qui \unsure{implique} la stabilité par facteurs directs, et \unsure{a fortiori|à fortiori} \unsure{aussi} la $=$ propriété pour \deleted{\unsure{$u^{*-1}(W)$}} $W'\overset{\mathrm{déf}}=(u^{*})^{-1}(W)$.}
\deleted{\unsure{donc}} \inserted{(et je} doute qu'il y \unsure{ait} \deleted{un triple de Quillen spécial} \inserted{des triples de Quillen spéciaux}
en dehors de ce cas), alors si
$(W,C,F)$ est un triple de Quillen
spécial, il en est de $=$ de $(W,\widetilde C=C^{\unsure{\natural|4}},F)$,
et je présume qu'on \deleted{\unsure{très}} aurait intérêt
à remplacer \unsure{le|ce} triple initial par celui-là.

\scanpage{111}
\mspage{98}
% archive sheet number: 110

6. \uline{Application à la construction des dérivateurs}
(associés aux triples de Quillen spéciaux).
% slanted
\margin{C'est le résultat \uline{principal} \unsure{de|dans} \unsure{ma|la} \unsure{partie} $\overline{\underline{\mathrm{XVII}}}$ des notes.}

\uline{Théorème} 6.1. Soit \deleted{\unsure{$M$}} $(W,C,F)$ un triple
de Quillen clos spécial \deleted{\unsure{dans}} \inserted{sur la} catégorie $M$
(cf déf.~5.6. p.~110, et cor.~5.7). Pour
tout $I$ dans Cat, considérons la catégorie
\[
 M(I)=\underline{\mathrm{Hom}}(I^{\unsure{o|\circ}},M),\tag*{(6.1.1)}
\]
et \deleted{posons} considérons
\[
 W(I),\deleted{\unsure{C}}F(I)\subset\mathrm{Fl}\,M(I)\tag*{\deleted{(6.1.2)}}
\]
définis par
\[
\begin{aligned}
 &u:X\mathrel{\unsure{-|\longrightarrow}}Y\ \text{dans }M(I)\text{ est dans }\unsure{W(I)|W(F)}\\
 &\text{(resp. }F(I))\ \text{ssi }\ \forall i\in\mathrm{Ob}\,I,\quad u(i):X(i)\longrightarrow Y(i)\\
 &\text{est dans }W\ \text{(resp. }F).
\end{aligned}\tag*{(6.1.2)}
\]
\deleted{Alors} \inserted{a)} Il existe une structure de
Quillen close (néc. unique, \inserted{\unsure{avec} $C(I)=$}) dans $M(I)$,
ayant $W(I)$ et $F(I)$ comme \deleted{équiv. fai} \inserted{\unsure{ensembles de|ens. des}}
quasi-isomorphismes et de fibrations
respectivement. On a donc aussi
\[
\left\{\begin{aligned}
 TF(I)&=F(I)\cap W(I)=\{u\in\mathrm{Fl}(M(I))\mid\\
 &\hspace{1.25em}u(i)\in TF\quad\forall i\in\mathrm{Ob}\,I\},\unsure{.}\\
 C(I)&=TF(I)^{*}.
\end{aligned}\right.\tag*{(6.1.3)}
\]
Cette structure de Quillen est \uline{spéciale}.

\scanpage{112}
\mspage{99}
% archive sheet number: 111

b) Soit $I_{0}$ la catégorie discrète ayant $=$
$=$ ensemble d'objets que $I$, et
\[
 \unsure{v|u}_{I}^{\deleted{\unsure{*|I}}}\ \text{ou }v:I_{0}\mathrel{\unsure{-|\longrightarrow}}I\tag*{(6.1.4)}
\]
le foncteur canonique. \unsure{Distinguons|Désignons} par
\deleted{\[
 C(\unsure{I|\,})\tag*{(6.1.5)}
\]}
\[
 C(I)=(TF(I))^{*},\qquad TC(I)=C(I)\cap W(I)=F(I)^{*}\tag*{(6.1.5)}
\]
les ensembles de cofibr. resp. de cofibr.
triviales dans $M(I)$, et \unsure{utilisons} les
$=$s notations $C(I_{0}),TC(I_{0})$ pour \unsure{Do.|Do}. On
a alors
\[
\begin{aligned}
 &\deleted{\unsure{C(I_{0})}}\ u:X_{0}\mathrel{-}Y_{0}\ \text{dans }M(I_{0})\\
 &\text{est dans } C_{0}\ \text{(resp. }TC_{0})\ \text{ssi pour}\\
 &\text{tt }i\ \text{dans }\unsure{i|I},\quad u(i):X_{0}(i)\mathrel{-}Y_{0}(i)\\
 &\text{est dans }C\ \text{(resp. }TC).
\end{aligned}\tag*{(6.1.6)}
\]
Enfin, on a
\[
\left\{\begin{aligned}
 C(I)&=(v_{!}(C(I_{0})))^{\sim},\\
 TC(I)&=(v_{!}(TC(I_{0})))^{\sim}.
\end{aligned}\right.\tag*{(6.1.7.)}
\]

c) Soit \deleted{\unsure{un}}
\[
 u:I\longrightarrow I'
\]
une flèche \deleted{\unsure{quel}} dans Cat, considérons
le couple de foncteurs adjoints
\diagram{\begin{tikzcd}[ampersand replacement=\&,column sep=large]
{(6.1.8)} \& M(I) \arrow[r,shift left=1.2ex,"{u_{!}}"] \& M(I') \arrow[l,shift left=1.2ex,"{u^{*}}"]
\end{tikzcd}}.

\scanpage{113}
\mspage{100}
% archive sheet number: 112

Ce couple satisfait aux conditions de Quillen
pour les couples de foncteurs adjoints entre
catégories de Quillen (Hom. Alg. Chap I,
\S~4, \unsure{th.|Th.} 3 p.~4.5), de sorte qu'ils
définissent par passage aux localisés un
couple de foncteurs adjoints
\diagram{\begin{tikzcd}[ampersand replacement=\&,column sep=large]
{(6.1.9)} \& \mathrm{Ho}_{\unsure{W|w}}(I) \arrow[r,shift left=1.25ex,"{u_{!}^{W}}"] \& \mathrm{Ho}_{W}(I') \arrow[l,shift left=1.25ex,"{u^{*}_{W}}"]
\end{tikzcd}}\unsure{,|.}
On a, plus précisément, les relations
\[
\left\{\begin{aligned}
 u_{!}(C(I))&\subset C(I')\\
 u_{!}(TC(I))&\subset TC(I')
\end{aligned}\right.\tag*{(6.1.10)}
\]
d'où
\[
 u_{!}(M(I)_{c})\subset M(I')_{c},\tag*{(6.1.11)}
\]
et de plus
\[
\begin{aligned}
 &\text{pour toute \unsure{flèches|flèche} }\inserted{f}\ \text{dans }M(I)_{c},\text{ telle}\\
 &\text{que }f\in W(I),\ \text{on a }\ u_{!}(f)\in W(I').
\end{aligned}\tag*{(6.1.12)}
\]
\inserted{i.e. $(u_{!})_{\unsure{c}}:M(I)_{c}\mathrel{-}M(I')_{c}$ est comp. avec les \unsure{localiseurs|localisations}}
D'autre part on a
\[
\left\{\begin{aligned}
 u^{*}(F(I'))&\subset F(I)\\
 u^{\unsure{*|4}}(TF(I'))&\subset TF(I)\\
 u^{*}(W(I'))&\subset W(I)
\end{aligned}\right.\tag*{(6.1.13)}
\]
donc
\[
\left\{\begin{aligned}
 &u^{*}(M(I')_{f})\subset M(I)_{f}\\
 &u^{*}_{f}:M(I')_{f}\mathrel{\unsure{-|\longrightarrow}}M(I)_{f}\ \text{est compatible}\\
 &\hspace{5em}\text{avec les localiseurs}.
\end{aligned}\right.\tag*{(6.1.14)}
\]

\scanpage{114}
\mspage{101}
% archive sheet number: 113

\unsure{d)} Supposons que $I$ soit \unsure{un|une} \deleted{\unsure{ensemble}} \deleted{\unsure{cat.}}
ordonné, et considérons l'ensemble
$C_{I^{\unsure{o|\circ}}}\unsure{\subset|\in}\mathrm{Fl}(M(I))$ défini \deleted{\unsure{dans}} \unsure{au} \S~2\textsuperscript{$*$}
% slanted
\margin{Au \S~2 c'était noté $C_{I^{o}}$ \unsure{ald|àld} $C_{I}$}
(cf page 25, cf aussi th.~2.1. p.~26, \deleted{\unsure{3}} et
ses corollaires.) On a alors
\[
\left\{\begin{aligned}
 C_{I^{o}}&\supset C(I)\\
 TC_{I}&\supset TC(I)
\end{aligned}\right.,\tag*{(6.1.15)}
\]
et on a égalité quand $I$ est \deleted{\unsure{la}} un
ens. ordonné \uline{artinien}\unsure{.|..}

\uline{Démonstration}. L'assertion a) est un
cas particulier de 5.8., appliqué au
couple de foncteurs adjoints
\diagram{\begin{tikzcd}[ampersand replacement=\&,column sep=large]
M(I_{0}) \arrow[r,shift left=1.2ex,"{\unsure{v|u}_{!}}"] \& M(I) \arrow[l,shift left=1.2ex,"{v^{*}}"]
\end{tikzcd}}.
Les conditions (5.8.2) sont satisfaites,
cela a été vu dans (4.8.). La prop.~5.8.
Nous dit donc l'existence d'une
structure de triple de Quillen close
sur $M(I)$, dont l'ens. des \deleted{équiv. faibles} \inserted{qu. iso.}
$W(I)$ et l'ens des fibrations $F(I)$, sont
donnés par la première formule (5.8.3)
et par la première formule (5.8.4),
ce qui équivaut aux définitions (6.1.2).
[NB Le fait que $M(I_{0})$, avec la
structure $W(I_{0}),C(I_{0}),$\inserted{$F(I_{0})$} évidente \inserted{(\unsure{produit|pr membres|pr. valeurs})}, est

\scanpage{115}
\mspage{102}
% archive sheet number: 114

\inserted{bien} une catégorie de Quillen spéciale close, est
$\pm$ tautologique. ].

\uline{Prouvons} b). La relation (6.1.6) est
tautologique. Les relations (6.1.7) sont
contenues ds 5.8. (cf $2^{\unsure{\mathrm{me}|\mathrm{e}}}$ relation
ds (5.8.3) pour $C'=\unsure{v|u}_{!}(C)^{\sim}$ ou
$C(I)=(v_{!}(C(I_{0})))^{\sim}$, et (5.8.5) pour
$TC(I)=(v_{!}(TC(I_{0})))^{\sim}$.

\uline{Prouvons} c). La première assertion
est, par Quillen, un cas particulier
des relations (6.1.10) à (6.1.14), c'est
elles qu'il s'agit d'établir. Les relations
(6.1.13) et par suite (6.1.14) sont
tautologiques en vertu des définitions,
\deleted{\unsure{donc}} il s'agit d'établir \uline{(6.1.10)}
et (6.1.12).

Pour les formules (6.1.10), j'utilise
le diagramme commutatif
\diagram{\begin{tikzcd}[ampersand replacement=\&,column sep=large]
I \arrow[d,dash,"{u}"'] \& I_{0} \arrow[l,"{v}"'] \arrow[d,"{u_{0}}"]\\
I' \& I'_{0} \arrow[l,"{v'}"']
\end{tikzcd}},
et je note que les foncteurs $u_{!}^{M},u_{0!}^{M}$, commutant aux petites limites inductives, \deleted{\unsure{comme}} \inserted{se}

\scanpage{116}
\mspage{103}
% archive sheet number: 115

\unsure{comportent|compatibles} bien \deleted{\unsure{par rapport}} \unsure{à} l'opération
\deleted{\unsure{tant aux|tendent à la} formation des ensembles}
\[
\Phi\mathrel{\unsure{\vdash\joinrel\dashv|\longmapsto}}\Phi^{\sim},
\]
quand $\Phi$ est une partie de $\mathrm{Fl}\deleted{\unsure{I}}M(I)$
resp de $\mathrm{Fl}\deleted{\unsure{I}}M(I_{0})$. \inserted{qui est $\unsure{Q|\mathcal{U}}$-\uline{accessible}, i.e.}
telle qu'il existe une \uline{petite} partie
$\Phi_{0}\subset\Phi$ \unsure{avec} $\Phi\subset\Phi_{0}^{\sim}$,
\inserted{plus précisément \unsure{---|...}}
$u_{!}(\Phi^{\sim})\subset(u_{!}(\Phi))^{\sim}$.

Ceci noté, on a par b
\[
 C(I)=(v_{!}(C(I_{0})))^{\sim}
\]
d'où
\[
 u_{!}(C(I))\subset
 \bigl(\underbrace{u_{!}\bigl(v_{!}(C(I_{0}))\bigr)}_{(uv)_{!}(C(I_{0}))}\bigr)^{\sim},
\]
et la deuxième $=$ s'écrit aussi
\[
 (v'_{!}u_{\deleted{\unsure{0}}0!}(C(I_{0})))^{\sim}
 \overset{?}{\subset}(v'_{!}(I'_{0}))^{\sim}=C(I').
\]
Donc tout revient à montrer que
\[
 u_{0!}(C(I_{0}))\subset C(I'_{0}),\tag*{(6.1.16)}
\]
ce qui résulte aussitôt du \unsure{calcul}
de $u_{0!}$ à coups de sommes directes,
et du fait que $C$ est stable par
sommes directes quelconques. Cela
prouve donc $u_{!}(C(I))\subset C(I')$, et on
prouve de $=$ $u_{!}(TC(I))\subset TC(I')$,
\deleted{\unsure{en se réduisant}} \inserted{qui ramène de $=$ au cas particulier} évident
\[
 u_{0!}(TC(I_{0}))\subset TC(I'_{0}).\tag*{(6.1.17)}
\]

